\section*{第 \ref{ch:b2:proba} 章 --- 可数空间上的概率}

\begin{solution}{exo:b2:proba:1}
由对称性，抽取次序诱导出一个
关于球 $1$ 与球 $2$ 的均匀随机相对次序，故 $\P(1 \text{ 先于 } 2) =
\frac12$。形式地说，交换抽取序列中球 $1$ 与球 $2$ 的位置，
是等可能结果之间的一个双射，
它把该事件与其补事件互换。在
球 $1, \dots, k$ 中，这 $k$ 个球的相对次序
在 $k!$ 种排列中均匀分布，其中球 $1$ 居首的有 $(k-1)!$ 种，
故概率为 $\frac{(k-1)!}{k!} = \frac1k$。
\end{solution}

\begin{solution}{exo:b2:proba:2}
$\frac{1}{k(k+1)} = \frac1k - \frac1{k+1}$，故 $\sum_{k\geq1} p_k$
裂项相消，和为 $1$，故为概率测度。偶数结果的概率为：
\[
\P(2\N^*) = \sum_{j=1}^{\infty}\frac{1}{2j(2j+1)}
= \sum_{j=1}^{\infty}\Bigl(\frac{1}{2j} - \frac{1}{2j+1}\Bigr)
= \frac12 - \frac13 + \frac14 - \frac15 + \cdots
\]
此即交错调和级数去掉首项
并变号：由于 $\ln 2 = 1 - \frac12 + \frac13 -
\frac14 + \cdots$（见\cref{ch:b2:series}），
\[
\P(2\N^*) = -\bigl(\ln 2 - 1\bigr) = 1 - \ln 2 \approx 0.307 .
\]
\end{solution}

\begin{solution}{exo:b2:proba:3}
设 $S$ 表示患病，$+$ 表示检验呈阳性。由
贝叶斯公式（\cref{thm:b2:proba:bayes}），取划分 $\{S, S^c\}$：
\[
\P(S \mid +)
= \frac{0.99 \times 10^{-4}}
       {0.99 \times 10^{-4} + 0.01 \times 0.9999}
= \frac{0.000099}{0.000099 + 0.009999}
\approx 0.0098 ,
\]
低于 $1\%$。尽管该检验号称“准确率为 99\%”，阳性
结果仍使你约有 $99\%$ 的概率身体健康：假阳性
在占绝大多数的健康人群中，其数量淹没了
极少数患病者中的真阳性。罕见病
的筛查检验结果，必须始终借助这一基率
计算加以解读。
\end{solution}

\begin{solution}{exo:b2:proba:4}
逐点考察 $\Omega$ 上的情形：$\omega \in \bigcup A_i$ 当且仅当某个因子
$1 - \mathbf{1}_{A_i}(\omega)$ 为零，于是
\[
\mathbf{1}_{\bigcup A_i}
= 1 - \prod_{i=1}^n\bigl(1 - \mathbf{1}_{A_i}\bigr)
= \sum_{\emptyset \neq J \subseteq \{1,\dots,n\}}
(-1)^{\abs J + 1}\prod_{i \in J}\mathbf{1}_{A_i} ,
\]
展开乘积并将 $1$ 移项。此时
$\prod_{i\in J}\mathbf{1}_{A_i} = \mathbf{1}_{\bigcap_{i \in J}
A_i}$，并以 $\P(\{\omega\})$ 为权重求和 ---
这一步合法：有限多个有界项，各族皆可和 ---
把每个示性函数化为对应事件的概率，得到
所求公式。
\end{solution}

\begin{solution}{exo:b2:proba:5}
设 $A_i$ 为事件“第 $i$ 封信在正确的信封中”。对于 $J$
其大小为 $k$，$\P\bigl(\bigcap_{i\in J}A_i\bigr) =
\frac{(n-k)!}{n!}$（固定 $k$ 封信，对其余信件作置换）。由
容斥原理，
\[
\P\Bigl(\bigcup A_i\Bigr)
= \sum_{k=1}^n (-1)^{k+1}\binom nk \frac{(n-k)!}{n!}
= \sum_{k=1}^n \frac{(-1)^{k+1}}{k!} ,
\]
于是
\[
\P(\text{无匹配})
= 1 - \P\Bigl(\bigcup A_i\Bigr)
= \sum_{k=0}^{n}\frac{(-1)^k}{k!}
\xrightarrow[n\to\infty]{} e^{-1} \approx 0.368 .
\]
恰有一个匹配：恰有一个不动点的排列，
由不动点的选取（$n$ 种方式）及
\emph{错排}（其余 $n - 1$ 封信的无匹配排列）唯一确定；
记 $D_{n-1} = (n-1)!\sum_{k=0}^{n-1}\frac{(-1)^k}{k!}$ 为
错排数（即前述和式乘以
$(n-1)!$），
\[
\P(\text{恰有一个匹配})
= \frac{n\,D_{n-1}}{n!}
= \frac{D_{n-1}}{(n-1)!}
= \sum_{k=0}^{n-1}\frac{(-1)^k}{k!}
\xrightarrow[n\to\infty]{} e^{-1} :
\]
极限情形下，“无匹配”与“恰有一个匹配”
同样可能，概率均为 $e^{-1}$。
\end{solution}

\begin{solution}{exo:b2:proba:6}
以首次投掷的结果为条件（乘法公式 /
\cref{thm:b2:proba:bayes}）：
\begin{itemize}
\item 首次掷出 T（概率 $1 - p$）：游戏
从头开始，持续超过 $n$ 次即等价于从该处起持续超过 $n - 1$
次，贡献为 $(1-p)\,q_{n-1}$；
\item 前两次掷出 HT（概率 $p(1-p)$）：掷两次后
重新开始，贡献为 $p(1-p)\,q_{n-2}$；
\item 前两次掷出 HH：游戏已经结束（$n$ 次之内，$n
\geq 2$）：贡献为 $0$。
\end{itemize}
因此 $q_n = (1-p)q_{n-1} + p(1-p)q_{n-2}$。其特征
方程 $r^2 = (1-p)r + p(1-p)$ 的根为
\[
r_\pm = \frac{(1-p) \pm \sqrt{(1-p)^2 + 4p(1-p)}}{2},
\]
且 $\abs{r_\pm} < 1$：事实上，多项式 $\chi(r) = r^2 -
(1-p)r - p(1-p)$ 满足 $\chi(1) = 1 - (1-p) - p(1-p) = p^2 >
0$ 且 $\chi(-1) = 1 + (1-p) - p(1-p) > 0$，而 $\chi(0) =
-p(1-p) < 0$：一根位于 $\intoo{-1}{0}$，另一根位于 $\intoo{0}{1}$。
于是 $q_n = \alpha r_+^n + \beta r_-^n \to 0$。事件“游戏
持续超过 $n$ 次”递减趋于“游戏永不结束”；由单调
连续性（见\cref{thm:b2:proba:continuity}）可得
$\P(\text{永不结束}) = \lim q_n = 0$：游戏几乎
必然结束。
\end{solution}

\begin{solution}{exo:b2:proba:7}
\emph{$\P(R_n) = 1/n$：}在前 $n$ 次抽取中，最后一次抽取的 $n$ 个
相对位置等可能（相对次序的均匀性），
$R_n$ 是“它为其中最大者”这一事件，
概率为 $1/n$。

\emph{无穷多个记录：}$\sum_n \P(R_n) = \sum 1/n =
\infty$，$R_n$ 相互独立（在本阶段不加证明地承认），因此
由博雷尔–坎泰利~2（\cref{thm:b2:proba:borelcantelli}）可得
$\P(\limsup R_n) = 1$：几乎必然地记录永不停止 --- 但
它们按对数的速度变得稀疏。

\emph{相邻记录：}由独立性，
\[
\sum_n \P(R_n \cap R_{n+1})
= \sum_n \frac{1}{n(n+1)} < \infty ,
\]
于是可用博雷尔–坎泰利~1：几乎必然地，只有有限多次
出现记录之后紧接着又是记录的情形。
引理的两部分相互配合：记录有无穷多个，
但（几乎必然地）最终不再出现两次连续的记录。
\end{solution}

\begin{solution}{exo:b2:proba:8}
\emph{第一部分：}$\sum \P(A_n) = \sum\frac{1}{n+1} = \infty$
在事件相互独立的情形下，由博雷尔–坎泰利~2 可得 $\P(\limsup A_n) = 1$。
每个 $A_n$ 发生的可能性越来越小，然而几乎每个
$\omega$ 都属于其中无穷多个。

\emph{不含独立性的反例：}取 $\Omega = \N^*$
以 $p_k = \frac{1}{k(k+1)}$ 为权重（见\cref{exo:b2:proba:2}），
并取 $A_n = \{k \in \N^* : k \geq n\}$。
于是
\[
\P(A_n) = \sum_{k \geq n}\Bigl(\frac1k - \frac1{k+1}\Bigr)
= \frac1n ,
\qquad
\sum_n \P(A_n) = \infty ,
\]
但 $A_n$ 是递减的，故 $\limsup_n A_n = \bigcap_n A_n =
\emptyset$，$\P(\limsup A_n) = 0$。仅凭 $\sum\P(A_n)$ 发散，
并不能说明什么：事件若堆积在越缩越小的
空间的一部分上 --- 独立性正是排除了这种
“巧合”。
\end{solution}

\begin{solution}{exo:b2:proba:9}
记 $q = 1 - p$，首次正面出现在第 $2j + 1$ 次
的概率为 $q^{2j}p$，于是
\[
\P(\text{奇数次}) = \sum_{j\geq0}q^{2j}p
= \frac{p}{1 - q^2} = \frac{1}{1 + q} .
\]
对于公平硬币：$\frac1{1 + 1/2} = \frac23$。（合理性检验：
奇数次应更可能，因为第 $1$ 次排在最前 ---
事实上 $\frac1{1+q} > \frac12$ 恒成立。）
\end{solution}

\begin{solution}{exo:b2:proba:10}
事件 $B_N = \bigcap_{n=1}^N A_n^c$ 递减趋于
$\bigcap_nA_n^c$，由余事件的独立性，
$\P(B_N) = \prod_{n=1}^N(1 - p_n)$；由单调
连续性（见\cref{thm:b2:proba:continuity}）可得上式中的极限。
取对数，$\prod(1 - p_n) > 0$ 当且仅当 $\sum-\ln(1 -
p_n) < \infty$。若 $\sum p_n < \infty$，则 $p_n \to 0$，且
$-\ln(1 - p_n) \sim p_n$，对数级数收敛。若
$\sum p_n = \infty$，则 $-\ln(1 - p_n) \geq p_n$ 迫使
其发散，于是乘积为 $0$。这与
博雷尔–坎泰利~2 相符：当 $\sum p_n = \infty$ 时，不仅
$\P(\text{没有 }A_n\text{ 发生}) = 0$，而且几乎必然地
有无穷多个 $A_n$ 发生。
\end{solution}

\begin{solution}{exo:b2:proba:11}
设箱子 $A$ 为首先被发现为空的那个，另一箱中装有
$k$ 个。这意味着：在前 $2n - k$ 次取用中，
恰有 $n$ 次取向 $A$、$n - k$ 次取向 $B$（顺序任意），
且第 $2n - k + 1$ 次取用又取向 $A$，
发现其为空。各次取用是独立的公平选择，
因此该事件的概率为
$\binom{2n-k}{n}2^{-(2n-k)}\cdot\frac12$；乘以二（空箱
可以是两个箱子中的任一个）得
\[
\P(\text{另一箱中有 }k) = \binom{2n-k}{n}\,2^{-(2n-k)} .
\]
当 $n = 1$ 时：$k = 1$ 给出 $\binom11 2^{-1} = \frac12$，而 $k
= 0$ 给出 $\binom21 2^{-2} = \frac12$，总和为 $1$，
这正如所应当的。
\end{solution}

\begin{solution}{exo:b2:proba:12}
(a) 若 $\P(\{n\}) = c$ 对所有 $n$ 成立，则由 $\sigma$-可加性
得 $1 = \sum_nc$，不可能：无论 $c = 0$（和为 $0$）
还是 $c > 0$（和为无穷）。$\N$ 上不存在
均匀概率。

(b) 若 $A \cap B = \emptyset$ 且 $d(A)$、$d(B)$ 存在，则
$\abs{(A \sqcup B)\cap\intint1n} = \abs{A\cap\intint1n} +
\abs{B\cap\intint1n}$，故 $d(A \sqcup B) = d(A) + d(B)$：
这说明此类集合对上具有有限可加性。每个单点集的计数
函数最终为常数，故其密度为 $0$，而 $d(\N^*) = 1$。
若 $d$ 是 $\sigma$-可加的，则 $\N^* =
\bigsqcup_k\{k\}$ 将给出 $1 = \sum_k 0 = 0$：密度只是
有限可加的，而非 $\sigma$-可加的 --- 该公理
确有实质内容。

(c) 令 $A = \bigcup_{k\geq0}\intint{4^k}{2\cdot4^k - 1}$
（区段自 $4^k$ 至 $2\cdot4^k - 1$）。当 $n = 2\cdot4^K -
1$ 时，计数为 $\sum_{k\leq K}4^k \sim \frac43 4^K$，从而
比值 $\to \frac23$；当 $n = 4^{K+1} - 1$ 时，计数
不变，比值 $\to \frac13$。该比值在
两极限 $\frac13$ 与 $\frac23$ 之间振荡，故不存在密度。
\end{solution}

\begin{solution}{pb:b2:proba:1}
\textbf{1.} 长度为 $n$ 的路径由其上升步的集合确定，
终点为 $k$ 意味着有 $u$ 个上升步和 $n - u$ 个下降步，
且 $u - (n - u) = k$，即 $u = \frac{n+k}2$：
当且仅当 $n + k$ 为偶数且 $\abs k \leq n$ 时可能，共有
$\binom{n}{(n+k)/2}$ 种方式。每条具体路径都是 $n$ 次投掷的公平乘积测度中的一点，
其概率为 $2^{-n}$。
因此 $\P(S_n = k) = N_n(k)2^{-n}$。

\textbf{2.} $S_n$ 与 $n$ 同奇偶，故 $S_{2n+1} \neq
0$；且 $u_n = N_{2n}(0)4^{-n} = \binom{2n}n4^{-n}$。取值为：
$u_1 = \frac12$，$u_2 = \frac6{16} = \frac38$，$u_3 =
\frac{20}{64} = \frac5{16}$。

\textbf{3.} $\dfrac{u_n}{u_{n-1}} =
\dfrac{\binom{2n}n}{4\binom{2n-2}{n-1}} =
\dfrac{(2n)(2n-1)}{4n^2} = \dfrac{2n-1}{2n} < 1$：递减。
由\cref{ex:b2:comparison:centralbinomial}，$\binom{2n}n \sim
\frac{4^n}{\sqrt{\pi n}}$，故 $u_n \sim \frac1{\sqrt{\pi n}}
\to 0$，而 $\sum u_n$ 通过与 $\sum
n^{-1/2}$ 比较而发散。

\textbf{4.} 设一条从 $1$ 到 $k$ 且触及 $0$ 的路径，
将其初始段（直到\emph{首次}到达
$0$）为止，关于水平轴作反射，所得为一条从
$-1$ 到 $k$ 的路径，且该操作是对合 --- 任一
从 $-1$ 到 $k \geq 1$ 的路径必经过 $0$，而将其初始段
再反射回去即可恢复原路径。因此触及的路径条数为
$N_{n-1}(k + 1)$（从 $-1$ 到 $k$ 的
位移为 $k + 1$）。从 $0$ 到 $k$ 的路径，使得 $>
0$（在时刻 $0$ 之后），首步为上升步，随后从
$1$ 到 $k$ 且在 $n - 1$ 步内不触及 $0$，共有
$N_{n-1}(k-1) - N_{n-1}(k+1)$ 条。

\textbf{5.} 记 $m = \frac{n+k}2$，利用
$\binom{n-1}{m-1} = \frac mn\binom nm$ 及 $\binom{n-1}{m} =
\frac{n-m}n\binom nm$：
\[
\frac{N_{n-1}(k-1) - N_{n-1}(k+1)}{N_n(k)}
= \frac{\binom{n-1}{m-1} - \binom{n-1}{m}}{\binom nm}
= \frac{m - (n - m)}{n} = \frac kn .
\]
当 $n = 3$，$k = 1$ 时：$N_3(1) = 3$ 条路径（$++-$，$+-+$，
$-++$），其中只有 $++-$ 保持为正（$+-+$ 回到
$0$，时刻为 $2$）：占三分之一，且 $\frac kn = \frac13$。

\textbf{6.} 由对称性，所求概率为 $2\P(S_i > 0\ \forall
i \leq 2n)$。对终点 $2k$ 求和，并利用
第 4 问（把 $n$ 换成 $2n$）：
\[
\P(S_i > 0\ \forall i) = 2^{-2n}\sum_{k\geq1}
\bigl(N_{2n-1}(2k-1) - N_{2n-1}(2k+1)\bigr)
= 2^{-2n}\,N_{2n-1}(1),
\]
这是一个裂项和。又 $N_{2n-1}(1) = \binom{2n-1}{n}$ 及
$2\binom{2n-1}n = \binom{2n}n$（帕斯卡公式），故所示
概率为 $2\cdot2^{-2n}\binom{2n-1}n =
\binom{2n}n4^{-n} = u_n$。

\textbf{7.} 事件 $D_n = \{S_i \neq 0,\ i \leq 2n\}$
递减，其交集为“永不返回”；由单调连续性
及第 6 问，$\P(\text{永不返回}) = \lim u_n =
0$：该随机游走几乎必然返回。此外 $f_n = \P(D_{n-1})
- \P(D_n) = u_{n-1} - u_n$，且由第 3 问
\[
u_{n-1} - u_n = u_n\Bigl(\frac{2n}{2n-1} - 1\Bigr) =
\frac{u_n}{2n-1};
\qquad
\sum_{n\geq1}f_n = u_0 - \lim u_n = 1 .
\]

\textbf{8.} $2n\,f_n = \frac{2n}{2n-1}u_n \geq u_n$，且
$\sum u_n = \infty$（第 3 问）：级数 $\sum 2nf_n$
发散。首次返回是必然的，但其平均
等待时间不是有限的 --- 该随机游走为 \emph{零常返}，
这一术语将在\cref{ch:b2:randomvar}中给出。

\textbf{9.} 事件“至少返回 $k$ 次”是
对 $0 < n_1 < \dots < n_k$ 所作的两两不交可数并，其中每一项为
事件“前 $k$ 次返回恰好发生在时刻
$2n_1, \dots, 2n_k$”。此类事件是
$k$ 个事件的交，它们依赖于互不相交的投掷块
$\intint1{2n_1}$、$\intint{2n_1+1}{2n_2}$、\dots，每一块
都需要一条新的随机游走，使其首次返回恰在
规定的步数之后发生；由各块的独立性，
其概率为 $f_{n_1}f_{n_2-n_1}\cdots
f_{n_k-n_{k-1}}$。按组求和（\cref{ch:b2:series}，
各项非负）：
\[
\P(\text{至少 }k\text{ 次返回})
= \Bigl(\sum_{n\geq1}f_n\Bigr)^{\!k} = 1^k = 1 .
\]
事件随 $k$ 递减，故由单调连续性，
$\P(\text{无穷多次返回}) = 1$：即常返性。

\textbf{10.} 由第 9 问，该随机游走无穷多次
离开 $0$ 的游历。每次游历的首步是
一次新的投掷，与之前的一切独立：
前 $m$ 次游历均向下出发的概率
为 $2^{-m}$。要到达 $1$，随机游走只需一次向上的游历起始
（自 $<0$ 出发者须先经过 $0$ 方能
到达 $1$，因步长为 $\pm1$），故 $\P(\text{从不击中 }1)
\leq 2^{-m}$ 对每个 $m$ 成立：该随机游走几乎必然击中 $1$。
按（几乎必然有限的）击中时刻进行分解，
该处重新开始的随机游走是从 $1$ 出发的新随机游走：由
归纳法，它几乎必然击中每个 $k \geq 1$，并由
对称性击中每个 $k \leq -1$。最后，从首次
访问 $k$ 时重新开始，第 9 问适用于新的随机游走：每个
格点几乎必然被无穷多次访问。

\textbf{11.} 事件 $A_n = \{S_{2n} = 0\}$ 远非
独立（处于 $0$ 的状态在时刻 $2n$ 出现，使得处于 $0$ 的
状态在时刻 $2n + 2$ 出现的可能性远大于 $u_{n+1}$），因此
博雷尔–坎泰利~2 不再适用，且事实上整个
第二部分的全部工作正是为了替代它。另一方向无需
独立性：\emph{若} $\sum\P(A_n)$ 收敛，
博雷尔–坎泰利~1 给出几乎必然只有有限次返回。
这一蕴涵是下文每个暂留性证明的
核心。

\textbf{12.} 时刻 $2n$ 的返回需要 $n$ 次向上与 $n$ 次
向下步：$\P(S_{2n} = 0) = \binom{2n}np^nq^n =
u_n(4pq)^n$，且 $4pq = 1 - (p - q)^2 < 1$，当 $p \neq
\frac12$。由 $u_n \leq 1$，级数 $\sum\P(S_{2n} = 0)$
被几何级数 $\sum(4pq)^n$ 所控制，故收敛。由
博雷尔–坎泰利~1，$\P(S_{2n} = 0 \text{ 无限次}) =
0$：几乎必然只有有限次返回。

\textbf{13.} 当 $n + k$ 为偶数时，$\P(S_n = k) =
\binom{n}{\frac{n+k}2}p^{\frac{n+k}2}q^{\frac{n-k}2}$；
二项式系数不超过中心项，且
$p^{\frac{n+k}2}q^{\frac{n-k}2} = (pq)^{n/2}(p/q)^{k/2}$，
由此得所述界 $\leq
2^n(pq)^{n/2}(p/q)^{k/2} = (4pq)^{n/2}(p/q)^{k/2}$，关于 $n$ 可和，
因 $\sqrt{4pq} < 1$。博雷尔–坎泰利~1：格点 $k$ 几乎必然
只被访问有限次；对 $k \in
\Z$ 的例外零概率事件之并仍为零概率（可数
次可加性）。几乎必然每个格点都只被访问有限次，
因此整数序列 $(S_n)$ 永久离开每个有界
窗口：$\abs{S_n} \to \infty$。

\textbf{14.} $\P(S_n \neq 0,\ 1 \leq n \leq 200) = u_{100} =
\binom{200}{100}4^{-100} \approx \frac1{\sqrt{100\pi}}
\approx 0.056$：$200$ 次公平抛掷从不打平的
概率超过二十分之一。$1/\sqrt{\pi n}$ 的衰减
极为缓慢：必然打平（第 7 问）
与极长的无打平区段并不矛盾 --- 初步
领略第四部分的反正弦现象。

\textbf{15.} 以第一步作条件。若 $X_1 = +1$，则
$T_1 = 1$，与 $f_1 = \frac12$ 相符。若 $X_1 = -1$，则
随机游走必须从 $-1$ 爬升到 $1$；由分块分解，
首次回到 $0$ 发生在时刻 $2n$ 的事件分解为：
先向下走一步，然后由 $-1$ 出发的新随机游走首次
到达 $0$ --- 等价地，新随机游走首次到达
$+1$ --- 历时 $2n - 1$ 步，或对称地向上。
两种符号的贡献相等：
\[
f_n = 2\cdot\tfrac12\,\P(T_1 = 2n - 1) = \P(T_1 = 2n-1) .
\]
因此 $\P(T_1 < \infty) = \sum f_n = 1$，而 $\sum_n(2n -
1)f_n = \sum_n u_n = \infty$（由第 7 问）：随机游走
几乎必然到达 $1$，平均时间却为无穷。

\textbf{16.} 将 $\{M_n \geq k\}$ 按终值
$S_n = m$ 作划分。当 $m \geq k$ 时，条件 $M_n \geq k$
自动成立。对 $m < k$，将路径在
\emph{首次}到达水平 $k$ 之后反射：这是
$\{M_n \geq k, S_n = m\}$ 与 $\{S_n = 2k - m\}$ 之间的双射（每条
终点为 $2k - m > k$ 的路径都访问 $k$；反射回去即为其
逆）。于是
\[
\P(M_n \geq k)
= \sum_{m > k}\P(S_n = m) + \P(S_n = k)
+ \sum_{m < k}\P(S_n = 2k - m)
= 2\P(S_n > k) + \P(S_n = k).
\]

\textbf{17.} 取偶数时刻 $2n$ 且 $k = 1$：$\P(S_{2n} = 1)
= 0$，且 $\P(S_{2n} > 1) = \P(S_{2n} \geq 2)$，故
\[
\P(M_{2n} \geq 1) = 2\P(S_{2n} \geq 2)
= \P(S_{2n} \geq 2) + \P(S_{2n} \leq -2)
= 1 - u_n .
\]
因此 $\P(S_i \leq 0\ \forall i \leq 2n) = u_n$：随机游走
前 $2n$ 步中从不领先，其概率恰与它
从不打平（第 6 问）的概率相同 --- 两个截然不同的事件，
却由同一个 $u_n$ 刻画。

\textbf{18.} $\{L_{2n} = 2k\} = \{S_{2k} = 0\} \cap
\{\text{投掷 } 2k+1, \dots, 2n \text{ 的随机游走
无零点}\}$。两事件依赖于互不相交的投掷段，
因此它们独立；第一个事件的概率为 $u_k$，第二个
事件的概率为 $u_{n-k}$，由第 6 问应用于新的
$(2n-2k)$ 步随机游走。于是 $\P(L_{2n} = 2k) = u_ku_{n-k}$。
由于 $L_{2n}$ 恰取值 $0, 2, \dots, 2n$，
这些概率之和为 $1$：
$\sum_{k=0}^nu_ku_{n-k} = 1$，这是一个二项恒等式，
由概率划分得到。

\textbf{19.} 对称性是显然的：$u_ku_{n-k} =
u_{n-k}u_k$。由于 $u_j$ 关于 $j$ 递减，乘积
$u_ku_{n-k}$ 对中心 $k$ 最小，在
端点 $k \in \{0, n\}$ 处最大，此时其值为 $u_n$；
定量地，$u_ku_{n-k} \approx \frac1{\pi\sqrt{k(n-k)}}$
在主体部分；与之相对，在两端 $u_n \approx \frac1{\sqrt{\pi n}}$，
两者形成对照。当 $n = 5$ 时：$\P(L_{10} = 0) = \P(L_{10} = 10) =
u_5 = \frac{63}{256} \approx 0.246$，而 $\P(L_{10} = 4) =
u_2u_3 = \frac38\cdot\frac5{16} = \frac{15}{128} \approx
0.117$。在一个长时间的公平博弈中，最后一次打平
最可能出现在最开始或最末尾附近：通常有一方
会在极长的区段内持续领先，
尽管硬币并无偏倚。

\textbf{20.} 图景如下：在时刻 $n$，随机游走处于尺度
$\sqrt n$ 量级（即第 3 问的二项散布 --- $u_n \sim
1/\sqrt{\pi n}$ 即中心峰的高度）；它
无限次返回 $0$，概率为 $1$（第二部分
），然而两次返回之间的等待时间的期望发散
（第 8 问），这正是单次游历能占据任一
时间跨度中的一个正比例；相应地，一场 $2n$ 步博弈的最后一次
打平分布开来，极端值
最有可能（第 18--19 问），而从不领先的
概率与从不打平的概率相同，同为缓慢衰减的 $u_n$
（第 17 问）。极限中的必然性，每一有限
时段内的持续性：这便是公平随机游走。

\textbf{21.} 对 $\{S_{2n} = 0\}$（$n \geq 1$）作划分，按
首次回返时刻 $2k$，$1 \leq k \leq n$：前 $2k$ 次抛掷
实现首次回返，其余 $2n - 2k$ 次抛掷实现一个新游走的
回返，各段相互独立：
$u_n = \sum_{k=1}^nf_ku_{n-k}$。两个级数
$U(x) = \sum u_nx^n$，$F(x) = \sum f_nx^n$ 的收敛半径均 $\geq
1$（系数属于 $\intcc01$），由柯西乘积
（\cref{ch:b2:powerseries}）可得，对 $0 \leq x < 1$，
\[
U(x) - 1 = \sum_{n\geq1}\Bigl(\sum_{k=1}^n
f_ku_{n-k}\Bigr)x^n = F(x)\,U(x),
\qquad\text{即}\qquad
U(x)\bigl(1 - F(x)\bigr) = 1 .
\]

\textbf{22.} 当 $x \uparrow 1$ 时，$U(x)$ 与 $F(x)$ 单调递增
（系数非负）；每个部分和 $\sum_{n\leq
N}u_n$ 都是 $\sum_{n\leq N}u_nx^n \leq U(x)$ 的极限，
$U(x) \uparrow \sum u_n \in \intoc0{+\infty}$，同理
$F(x) \uparrow f = \sum f_n$。若 $\sum u_n = \infty$：$1 -
F(x) = 1/U(x) \to 0$，故 $f = 1$。若 $\sum u_n = S < \infty$：
$1 - f = 1/S > 0$，故 $f < 1$。检验：公平游走，$\sum u_n =
\infty$ 且 $f = 1$（第 3、7 题）；有偏游走，$\sum
u_n(4pq)^n < \infty$，相应地 $f = 1 -
1/\sum_{n\geq0}u_n(4pq)^n < 1$，与
回返次数几乎必然有限这一事实相一致（第
12 题）。

\textbf{23.} 对游走的四个步 $(\pm1, 0), (0, \pm1)$，
$\Z^2$ 游走中，$U = X + Y$ 与 $V = X - Y$ 的增量为：
$(+,+)$ 对应 $(1,0)$，$(+,-)$ 对应 $(0,1)$，$(-,+)$ 对应
$(0,-1)$，$(-,-)$ 对应 $(-1,0)$ --- 各组符号对的概率均为
$\frac14 = \frac12\cdot\frac12$，因此两个坐标游走
$(U_n)$ 与 $(V_n)$ 为相互独立的
$\Z$ 上公平游走。由于 $S^{(2)}_{2n} = (0,0)$ 当且仅当 $U_{2n} = 0$
且 $V_{2n} = 0$，
\[
\P\bigl(S^{(2)}_{2n} = (0,0)\bigr) = u_n^2 \sim \frac1{\pi
n}, \qquad \sum_nu_n^2 = \infty .
\]
第 21 题的更新恒等式与第 22 题的二分性
均未用到一维性（仅用到
按首次回返的分解及不相交分块的
独立性），因此由 $\sum u_n^{(2)} = \infty$ 得 $f^{(2)} =
1$，并借助第 9 题的论证加强此结论：
$\Z^2$ 上的游走几乎必然无限次回返原点。

\textbf{24.} 利用在本阶段不加证明地承认的界 $\P(S^{(3)}_{2n} = 0)
\leq Cn^{-3/2}$，级数收敛，由博雷尔--坎泰利~1
可知几乎必然只有有限次回返：
$\Z^3$ 上的游走是暂留的（同一界以指数 $-d/2$ 处理
每个 $d \geq 3$）。综上：\emph{波利亚
定理} --- 简单随机游走在 $\Z$ 与
$\Z^2$ 上常返，在 $\Z^d$（$d \geq 3$）上暂留。醉汉能找到回家的路；
醉鸟则未必。

\textbf{25.} 路径计数与反射给出了精确的
规律（$u_n$、选票定理、$f_n$、最大值、最后
一次零点）；单调连续性将每一个极限陈述
（“至少回返一次”、“无穷多次”）
转化为有限时域概率的极限；
不相交分块的独立性为更新
分解提供了动力（第 9、18、21 题） --- 它是
马尔可夫性质的可数骨架；博雷尔--坎泰利~1 是
暂留性的利器（第 12--13、24 题），无需
独立性；更新恒等式将一切归结为
二分性 $\sum u_n = \infty \iff$ 常返。唯一的
解析输入是局部估计 $u_n \sim 1/\sqrt{\pi
n}$：其平方 $1/(\pi n)$ 仍然发散（维数 $2$，
常返），而 $n^{-3/2}$ 收敛（维数 $3$，
暂留） --- 波利亚定理归根结底是
关于 $\sum n^{-d/2}$ 发散性的命题。
\end{solution}
